\section{Datasets}
\label{sec:datasets}

The corpus retains source and target identity for the paired audit.
All inputs are aligned to July~1, 2026 (UTC): network-path observations
are combined with IP and AS annotations and cable lifecycle records for
the analysis date. The following datasets supply the path population,
the annotations used in projection, and the geographic grouping used
in the comparisons.

\subsection{RIPE Atlas Measurements}
\label{subsec:atlas_corpus}

Our network observations comprise 18 public RIPE Atlas IPv4 traceroute
measurements collected from 00:00 to 01:00 UTC. Thirteen measurements cover
independently operated anycast DNS Root services. Wikipedia, Reddit, and
\texttt{assets.nflxext.com} provide three application-facing observations.
Measurements 5051 and 5151 use dynamic targets as multi-target UDP and ICMP
topology references. Table~\ref{tab:atlas_measurements} summarizes the corpus
after canonical trace-identity deduplication.

\begin{table*}[ht]
  \centering
  \caption{RIPE Atlas measurement corpus.}
  \label{tab:atlas_measurements}
  \footnotesize
  \setlength{\tabcolsep}{4.2pt}
  \begin{tabular}{p{0.15\textwidth}p{0.27\textwidth}p{0.30\textwidth}rr}
    \toprule
    Group & Target & MSM ID(s) & Raw & Valid \\
    \midrule
    DNS Roots & A-M Root
      & 5009, 5010, 5011, 5012, 5013, 5004, 5014, 5015, 5005, 5016,
        5001, 5008, 5006
      & 534,276 & 317,873 \\
    Applications & Wikipedia; Reddit; Netflix Assets
      & 86710103; 176906957; 176517335
      & 115,596 & 115,596 \\
    Topology references & IPv4 UDP; IPv4 ICMP
      & 5051; 5151
      & 57,442 & 57,442 \\
    \midrule
    \textbf{Total} & \textbf{18 measurements} & -
      & \textbf{707,314} & \textbf{490,911} \\
    \bottomrule
  \end{tabular}
\end{table*}

Each record retains the destination address returned by RIPE Atlas. This
preserves the multi-target structure of the two dynamic-target measurements
and the destination instances selected by service measurements during the
aligned window.

The two dynamic-target measurements provide a broader-target reference
for the service populations. They sample many destinations during the
window, whereas service measurements retain the destinations selected for
each service. The comparison describes these observed populations;
their different target-selection processes do not constitute a controlled
test of service deployment or target-selection effects.

\subsection{Supporting Datasets}
\label{subsec:supporting_data}

Table~\ref{tab:supporting_datasets} lists the supporting datasets and versions.
RIPE Atlas probe metadata supplies the probe country and source-network
context~\cite{ripe_atlas}. IPinfo Location and ASN databases annotate visible
hops with geographic and AS information~\cite{ipinfo}.

The Submarine Cable Map provides the cable and landing-point inventory,
including lifecycle and organizational fields
~\cite{telegeography_submarine_map}. The analysis includes systems in service
on the measurement date. CAIDA AS Relationships and the derived owner-AS
associations supply supplementary ranking context
~\cite{caida_as_relationships}. A versioned country taxonomy classifies probe
countries as island or archipelagic, landlocked, or other coastal.

\begin{table*}[htb]
  \centering
  \caption{Supporting datasets.}
  \label{tab:supporting_datasets}
  \footnotesize
  \begin{tabular}{p{0.30\textwidth}p{0.25\textwidth}p{0.35\textwidth}}
    \toprule
    Dataset & Version & Content \\
    \midrule
    RIPE Atlas probe metadata & 2026-07-01 & Probe country and source ASN \\
    IPinfo Location MMDB & 2026-07-01 & Hop geolocation \\
    IPinfo ASN MMDB & 2026-07-01 & IP-AS mapping \\
    Submarine Cable Map & 2026-07-01 & 694 cables; 1,916 landing points \\
    CAIDA AS Relationships & 2026-07-01 release & AS relationships \\
    Owner-AS mapping & 2026-07-01 inputs & Owner-AS associations \\
    Country geography taxonomy & 2026-07-01 & Operational geography classes \\
    \bottomrule
  \end{tabular}
\end{table*}
